\documentclass[a4paper,12pt,oneside]{book} % Formato documento: foglio a4, font size 12, solo fronte
\usepackage[utf8]{inputenc} % Codifica caratteri
\usepackage[lmargin=2cm,rmargin=2cm,tmargin=2.5cm,bmargin=2cm]{geometry} % Bordi pagina come da formato tesi Unibs
\linespread{1.5} % Interlinea
\setlength{\parindent}{0pt} % No spazio prima dei paragrafi
\setlength{\parskip}{1em}
\setlength{\headheight}{15pt}
\addtolength{\topmargin}{-2.5pt}

\usepackage[backend=biber,sorting=none]{biblatex} % Per poter scrivere la bibliografia su un file a parte
\addbibresource{refs.bib} % File contentente la bibliografia

\usepackage{titlesec} % Per modificare il formato dei capitoli
\titleformat{\chapter}[hang]{\normalfont\Huge\bfseries}{\thechapter.}{8pt}{\Huge\bfseries}
\titlespacing*{\chapter}{0pt}{-40pt}{20pt}

\usepackage{fancyhdr} % Per modificare l'intestazione delle pagine
\pagestyle{fancy} % Customizzazione testata e piè di pagina
\renewcommand{\chaptermark}[1]{\markboth{#1}{#1}}
\fancyhf{}
\fancyhead[L]{\thechapter. \leftmark}
\fancyfoot[C]{\thepage}

\usepackage{graphicx} % Per poter inserire immagini
\usepackage{float} % Per controllare il posizionamento di immagini e tabelle
\usepackage{amsmath} % Per inserire equazioni
\usepackage[hidelinks]{hyperref} % Per inserire URL
\usepackage{xurl} % Per consentire agli URL di spezzarsi a capo
\usepackage{enumitem}
\usepackage{siunitx} % Necessario per usare \setitemize
\setitemize{noitemsep,topsep=0pt,parsep=0pt,partopsep=0pt} % Per rendere le liste più compatte

\title{Implementation of The Multiple Knapsack Problem in the Kernel Search}
\author{Matteo Beatrice, Luca Cotti, Giacomo Golino}
\date{December 2020}

\begin{document}
    \begin{titlepage}
    \begin{figure}[H] % Logo Unibs
        \centering
        \includegraphics[width=72.4mm]{images/logo_unibs}
    \end{figure}

    \begin{center}
        \LARGE{\uppercase{Università degli Studi di Brescia}}\\ % Doppio backslash = vai a capo
        \vspace{5mm} % Spazio verticale
        \large{\uppercase{Dipartimento di Ingegneria dell'informazione}}\\
        \vspace{5mm}
        \large{Corso di Laurea Magistrale in Ingegneria Informatica}\\
    \end{center}

    \vspace{10mm}

    \begin{center}
        \LARGE{\textbf{Implementation of The Multiple Knapsack Problem in the Kernel Search}}\\
    \end{center}

    \vspace{10mm}

    \begin{flushleft}
        \large
        \textbf{Docente:}\\
        Ch.ma Prof.sa Renata Mansini
    \end{flushleft}

    \begin{flushright}
        \large
        \textbf{Studenti:}\\
        Matteo Beatrice (739848)\\
        Luca Cotti (719204)\\
        Giacomo Golino (719210)
    \end{flushright}

    \vspace*{\fill} % Spazio bianco per mettere il testo in fondo alla pagina

    \rule{0.8\textwidth}{0.6pt}\\ % Linea nera
    \centering{\Large{Anno Accademico 2021/2022}}
\end{titlepage} %Frontespizio

    \tableofcontents
    \hyphenpenalty=3000 % Penalizza lo spezzamento delle parole a capo

    % I capitoli veri e propri sono su pagine a parte
    \chapter{Introduction}

In this report we describe the process of adapting the Kernel Search,
a generic heuristic framework that can be successfully used on a variety of
problems, to the specific model of the Multiple Knapsack Problem.\\
While Kernel Search provides good solutions in a short amount of time,
tuning for a specifi problem becomes necessary to improve the performance of the algorithm
and the quality of the solution found.

In the first chapter, the Multiple Knapsack Problem is introduced,
along with a few real life examples to better explain the model.\\
Chapter~\ref{ch:kernel-search} contains a brief description of the Kernel Search algorithm,
along with a few basic improvements to the default algorithm.\\
In chapter~\ref{ch:implementation} we describe the technical
specifications of our implementation of the Kernel Search algorithm,
while in chapter~\ref{ch:computational-experiments} we study
the performance of the implementation and the test instances used as benchmark.\\
In chapter~\ref{ch:improvements} we detail all
the methods that we have introduced to better adapt the Kernel Search
algorithm to the Multiple Knapsack Problem.\\
Finally, in chapter~\ref{ch:testing} we report the results of the testing we have
executed on the modified Kernel Search algorithm.

    \chapter{The Multiple Knapsack Problem}
Given a set of \(m\) knapsacks with a known \(\text{capacity } c_{i}\) \((i=1,\dots,m)\) and a set of \(n\) items with
known \(\text{profit } p_{j}\) and \(\text{weight } w_{j}\) \((j=1,\dots,n)\) the Multiple Knapsack Problem (MKP) consists in
selecting \(m\) disjoint subsets of items (one subset per knapsack) such that the total weight of the items in the knapsack
does not exceed its capacity, and the profit of the selected items is maximized.

\section{Mathematical model}
There are quite a few formulations for the MKP, with different degrees of
complexity and performance.
The formulation that was used in this project is the classical and most intuitive
one, that uses the following binary variables:
\[x_{ij}=
\begin{cases}
    1 & \quad \text{if item } j \text{ is packed into knapsack } i\\
    0 & \quad \text{otherwise}
\end{cases}
\]

The model is as follows:
\begin{equation}
    \label{eq:obj}
    \max{z} = \sum_{i=1}^{n} \sum_{j=1}^{n} p_{j} x_{ij}
\end{equation}
\begin{equation}
    \label{eq:cap}
    \sum_{j=1}^{n} w_{j} x_{ij} \leq c_{i} \qquad i=1,\dots,m
\end{equation}
\begin{equation}
    \label{eq:sel}
    \sum_{i=1}^{m} x_{ij} \leq 1 \qquad j=1,\dots,n
\end{equation}
\begin{equation}
    x_{ij} \in \{0,1\} \qquad i=1,\dots,m \quad j=1,\dots,n
\end{equation}

The objective function~\eqref{eq:obj} maximises the profit of the items in the knapsacks.\\
Constraints~\eqref{eq:cap} impose that the capacity of each knapsack is respected, while constraints~\eqref{eq:sel}
ensure that each item is packed in no more than one knapsack.

\section{Real life examples}
The MKP has various real life examples, in various fields.

The most straightforward example might be the cargo organization for a transport company.
In this scenario, the knapsack may be the trucks or containers.
Their capacity is the effective capacity of the truck/container.
The items are the objects that must be delivered.
Their weight is the actual weight of the item, while the profit is the priority of the delivery.

Another example might be the project selection:
The knapsacks may be the years of development of the project, and their capacity is the available budget for
that year.
The items can be the available projects: the weight is the cost of the project, the profit is the monetary gains
from the project completion.
    \chapter{Implementation and Test Instances}


\section{Kernel Search}
The code provided during the course implements a generic kernel search
which uses GUROBI as solver, and has the following default paramenters:
\begin{itemize}
    \item Presolve = 2 (aggressive presolve)
    \item Time limit = 30 seconds
    \item Threads = 12
    \item MIPGap = \num{1e-12}
    \item Sort variables by profit and reduced cost
    \item Kernel size = 15\% of the number of variables (rounded to the nearest integer)
    \item Bucket size = 2.5\% of the number of variables (rounded to the nearest integer)
    \item Number of iterations = 2
    \item Time limit for each bucket = 5 seconds
\end{itemize}
The code was initially modified to implement the MKP model, and it was then
refactored to improve its quality and efficiency.
The algorithm itself and the parameters were not changed at this time.\\
The code along with extensive documentation can be found at
\url{https://github.com/Golino98/KernelSearchGolinoCottiBeatrice}.

\section{Test Instances}
The test instances for the MKP were kindly provided to us by the authors of~\cite{mkp:2019}.\\
There are in total 2100 test instances, organized in
five directories of increasing complexity:
\begin{itemize}
    \item SMALL
    \item FK\_1
    \item FK\_2
    \item FK\_3
    \item FK\_4
\end{itemize}
SMALL contains 180 instances, while FK\_1, FK\_2, FK\_3, FK\_4 contain 480 instances each.
The instances in the first three directories were already solved to the optimum.
Of those in the last two folders only a few were already solved: the unsolved
ones only had an upper and lower bound.
\newpage
\section{How to read an Instance}
Every instance was generated randomly by the authors of~\cite{mkp:2019}.
The name of each instance contains different information, in order to fastly understand the dimension of the problem represented.
The first integer represents the knapsacks number, while the second indicates then number of the items.
The last integer is used to differentiate the instances belonging to the same class.
Our test instances are text files, in which are contained the data related to the knapsacks and their \textbf{capacity} and items with their \textbf{profit} and \textbf{weight}. The structure of each of them shown below:

	\begin{flushleft} 
		$\begin{rcases*}
			3 \quad
		\end{rcases*} \rightarrow \textit{\textbf{m} knapsacks.}$
	\end{flushleft}
	\begin{flushleft}
		$\begin{rcases*}
			7 \quad
		\end{rcases*} \rightarrow \textit{\textbf{n} items.}$
	\end{flushleft}

	\begin{flushleft}
		$\begin{rcases*}
			908 \quad
			\\834 
			\\675 
		\end{rcases*} \rightarrow \textit{\textbf{capacities} of the m knapsacks.}$
	\end{flushleft}
	\begin{flushleft}
		$\begin{rcases*}
			264	\quad 430 \quad
			\\606 \quad	945
			\\268 \quad	409
			\\619 \quad	591
			\\958 \quad	839
			\\972 \quad	818
			\\723 \quad	71
		\end{rcases*} \rightarrow \textit{\textbf{weight} and \textbf{profit} of the n items.}$
	\end{flushleft}

For most of the instances without the exact solution, we were able to
find the exact using an exact GUROBI solver, with the following configuration:
\newpage
\begin{itemize}
    \item Presolve = 2 (aggressive presolve)
    \item Time limit = 90 hours
    \item Threads = 12
    \item MIPGap = \num{1e-12}
\end{itemize}

The code can be found at \url{https://github.com/Golino98/ProgettoTestingPorzioniCodiceKernelSearch}

To simplify the testing of the changes to the kernel search algorithm,
we only considered 40 instances: 6 for each directory,
plus 10 that couldn't be solved to the optimum.

Tables~\ref{tab:exact_opt} and~\ref{tab:exact_noopt} represent,
respectively, the optimum values of the 30 instances solved to
the optimum, and the upper and lower bound for the 10
that couldn't be solved.

\begin{table}[]
    \centering
    \resizebox{0.6\columnwidth}{!}{
        \begin{tabular}{@{}lll@{}}
            \toprule
            Directory              & Instance                                   & OPT    \\ \midrule
            \multirow{6}{*}{FK\_1} & random10\_60\_1\_1000\_1\_12               & 23064  \\
            & random10\_100\_2\_1000\_1\_10              & 29800  \\
            & random12\_48\_3\_1000\_1\_16               & 11865  \\
            & random15\_45\_1\_1000\_1\_13               & 15160  \\
            & random15\_75\_3\_1000\_1\_16               & 18321  \\
            & random30\_60\_4\_1000\_1\_14               & 11017  \\ \midrule
            \multirow{6}{*}{FK\_2} & random20\_120\_1\_1000\_1\_12              & 47823  \\
            & random20\_200\_2\_1000\_1\_14              & 53618  \\
            & random24\_96\_3\_1000\_1\_18               & 23912  \\
            & random30\_90\_4\_1000\_1\_20               & 26038  \\
            & random30\_150\_2\_1000\_1\_17              & 44518  \\
            & random60\_120\_4\_1000\_1\_19              & 20463  \\ \midrule
            \multirow{6}{*}{FK\_3} & random30\_180\_1\_1000\_1\_20              & 75618  \\
            & random30\_300\_2\_1000\_1\_4               & 85147  \\
            & random36\_144\_3\_1000\_1\_17              & 39245  \\
            & random45\_135\_4\_1000\_1\_16              & 37005  \\
            & random45\_225\_3\_1000\_1\_15              & 56199  \\
            & random90\_180\_4\_1000\_1\_19              & 30047  \\ \midrule
            \multirow{6}{*}{FK\_4} & random50\_300\_1\_1000\_1\_16              & 119973 \\
            & random50\_500\_2\_1000\_1\_14              & 135583 \\
            & random60\_240\_3\_1000\_1\_14              & 61180  \\
            & random75\_225\_4\_1000\_1\_17              & 60458  \\
            & random75\_375\_2\_1000\_1\_5               & 101917 \\
            & random150\_300\_2\_1000\_1\_12             & 55707  \\ \midrule
            \multirow{6}{*}{SMALL} & probT1\_0U\_R50\_T002\_M010\_N0040\_seed06 & 16510  \\
            & probT1\_0U\_R50\_T002\_M020\_N0060\_seed01 & 26593  \\
            & probT1\_0U\_R50\_T002\_M010\_N0040\_seed05 & 15534  \\
            & probT1\_1W\_R50\_T002\_M010\_N0040\_seed09 & 12724  \\
            & probT1\_1W\_R50\_T002\_M010\_N0060\_seed10 & 19652  \\
            & probT1\_1W\_R50\_T002\_M020\_N0020\_seed10 & 3429   \\
            & probT1\_1W\_R50\_T002\_M020\_N0040\_seed10 & 12405  \\ \bottomrule
        \end{tabular}
    }
    \caption{Exact solutions for the 30 instances already solved to the optimum}
    \label{tab:exact_opt}
\end{table}

\begin{table}[!htbp]
    \centering
    \begin{tabular}{@{}llll@{}}
        \toprule
        Directory              & Instance                      & LB    & UB    \\ \midrule
        \multirow{5}{*}{FK\_3} & random30\_180\_2\_1000\_1\_2  & 47911 & 47966 \\
        & random36\_144\_2\_1000\_1\_20 & 42988 & 43004 \\
        & random45\_135\_3\_1000\_1\_12 & 37680 & 37681 \\
        & random45\_135\_4\_1000\_1\_14 & 52666 & 33756 \\
        & random45\_225\_2\_1000\_1\_20 & 68070 & 68125 \\ \midrule
        \multirow{5}{*}{FK\_4} & random50\_300\_2\_1000\_1\_5  & 80459 & 80483 \\
        & random60\_240\_1\_1000\_1\_17 & 95602 & 95792 \\
        & random60\_240\_2\_1000\_1\_15 & 66497 & 66529 \\
        & random60\_240\_3\_1000\_1\_18 & 59637 & 59648 \\
        & random75\_225\_4\_1000\_1\_20 & 58338 & 58340 \\ \bottomrule
    \end{tabular}
    \caption{Upper and lower bounds for the 10 instances not solved to the optimum}
    \label{tab:exact_noopt}
\end{table}



\section{Performance of the default kernel search}


\section{Performance of random variable sorting}

    \printbibliography
\end{document}